% ==============================================================================
% Generic thesis / dissertation template (skeleton)
% Uses the standard `book` class (in every TeX Live install; book provides
% \frontmatter / \mainmatter / \backmatter and \chapter, which report lacks)
% so this file compiles out-of-the-box. Many universities ship their own .cls (e.g.
% mit-thesis, stanford-thesis, Cambridge, TsinghuaThuThesis, UCASthesis);
% for an actual submission download the institution's class and replace the
% \documentclass line below.
%
% Conventions:
%   Layout     : single-sided review copy (oneside) or print-ready (twoside)
%   Body font  : Times Roman 12pt (mathptmx)
%   Math font  : Times (mathptmx)
%   Line spacing: 1.5 for body text (per most graduate-school style guides)
%   Margin     : 1in (US letter) / 25 mm (A4) on all sides, check binding offset
%   Citation   : numeric [1] (or author-year per institution)
%   Headings   : \chapter / \section / \subsection, numbered
% Frontmatter : title page, committee page, abstract, acknowledgements,
%               dedication, TOC, list of figures, list of tables, notation
% Main        : Introduction, Literature Review, Methods, Results, Conclusion
% Backmatter  : bibliography, appendices
% ==============================================================================
\documentclass[12pt,a4paper,oneside]{book}
% For a print-ready book:
% \documentclass[12pt,a4paper,twoside,openright]{book}

\usepackage[T1]{fontenc}
\usepackage[utf8]{inputenc}
\usepackage{mathptmx}            % Times text + math
\usepackage{helvet}
\usepackage{courier}
\usepackage[a4paper,top=30mm,bottom=30mm,left=35mm,right=25mm]{geometry}
% For US letter use:
% \usepackage[letterpaper,top=1in,bottom=1in,left=1.5in,right=1in]{geometry}

\usepackage{amsmath,amssymb,amsfonts,amsthm}
\usepackage{graphicx}
\usepackage{booktabs}
\usepackage{array}
\usepackage{xcolor}
\usepackage{setspace}
\usepackage{fancyhdr}
\usepackage{tocbibind}           % add bib/appendix entries to the TOC
\usepackage{titlesec}
\usepackage{algorithm}
\usepackage{algpseudocode}
\usepackage{acronym}             % first-use expansion of abbreviations
\usepackage{lipsum}              % placeholder text (delete before submission)

% ---- 1.5 line spacing for the body (most graduate-school guides require it) ----
\onehalfspacing

% ---- Headers / footers ----
\pagestyle{fancy}
\fancyhf{}
\fancyhead[L]{\nouppercase{\leftmark}}
\fancyhead[R]{\thepage}
\renewcommand{\headrulewidth}{0.4pt}

% ---- Section heading tweaks (matches most thesis guides) ----
\titleformat{\chapter}[display]
  {\normalfont\huge\bfseries}{\chaptername\ \thechapter}{20pt}{\Huge}
\titlespacing*{\chapter}{0pt}{-20pt}{40pt}

% ---- Citations: numbered ----
\usepackage[numbers,sort&compress]{natbib}
\bibliographystyle{plainnat}     % numbered, alphabetical; switch to unsrtnat for citation-order
% For author-year use:
% \usepackage[authoryear,round]{natbib}
% \bibliographystyle{plainnat}

% \usepackage{hyperref}           % load last if used
% \usepackage{bookmark}

% ==============================================================================
% Frontmatter placeholders (edit each \newcommand below)
% ==============================================================================
\newcommand{\ThesisTitle}{A Concise and Informative Title of the Thesis}
\newcommand{\ThesisAuthor}{Author Full Name}
\newcommand{\ThesisDegree}{Doctor of Philosophy}
\newcommand{\ThesisDepartment}{Department of X}
\newcommand{\ThesisInstitution}{University of Y}
\newcommand{\ThesisYear}{20YY}
\newcommand{\ThesisAdvisor}{Advisor Full Name}
% Committee members (add or remove as needed):
\newcommand{\CommitteeChair}{Advisor Full Name, Chair}
\newcommand{\CommitteeMemberA}{Committee Member A}
\newcommand{\CommitteeMemberB}{Committee Member B}
\newcommand{\CommitteeMemberC}{Committee Member C, External Examiner}

\begin{document}

% ==============================================================================
% Front matter
% ==============================================================================
\frontmatter
\pagenumbering{roman}

% ---- Title page ----
\begin{titlepage}
\centering
\vspace*{2cm}
{\Huge \ThesisTitle \par}
\vspace{2cm}
{\Large \ThesisAuthor \par}
\vspace{1cm}
A dissertation submitted in partial fulfilment of\\
the requirements for the degree of\\
\ThesisDegree
\vfill
\ThesisDepartment\\
\ThesisInstitution\\
\ThesisYear
\end{titlepage}

% ---- Committee / examiner page ----
\chapter*{Committee}
\noindent This thesis has been examined and approved by the following committee:
\vspace{1em}

\noindent\textbf{\CommitteeChair}\\
\ThesisDepartment, \ThesisInstitution

\vspace{0.6em}
\noindent\textbf{\CommitteeMemberA}\\
\ThesisDepartment, \ThesisInstitution

\vspace{0.6em}
\noindent\textbf{\CommitteeMemberB}\\
\ThesisDepartment, \ThesisInstitution

\vspace{0.6em}
\noindent\textbf{\CommitteeMemberC}\\
Department of Z, External Institution

% ---- Abstract ----
\chapter*{Abstract}
\addcontentsline{toc}{chapter}{Abstract}
% TODO: 250-350 words. State the research question, the approach, the key
% contributions, and the main findings. Many guides allow up to 350 words.

% ---- Acknowledgements ----
\chapter*{Acknowledgements}
\addcontentsline{toc}{chapter}{Acknowledgements}
% TODO: thank the advisor, committee, funding, collaborators, family.

% ---- Dedication (optional) ----
% \chapter*{Dedication}
% \addcontentsline{toc}{chapter}{Dedication}
% \begin{center}\textit{To ...}\end{center}

% ---- Notation / abbreviations (optional but common in STEM theses) ----
\chapter*{List of Abbreviations}
\addcontentsline{toc}{chapter}{List of Abbreviations}
% TODO: use the `acronym` package, e.g.
% \begin{acronym}[TDMA]
%   \acro{CNN}{Convolutional Neural Network}
%   \acro{DNN}{Deep Neural Network}
% \end{acronym}

% ---- Table of contents, list of figures, list of tables ----
\tableofcontents
\listoffigures
\listoftables

% ==============================================================================
% Main matter
% ==============================================================================
\mainmatter
\pagenumbering{arabic}

% ---- Chapter 1: Introduction ----
\chapter{Introduction}
\label{ch:intro}
% TODO: research motivation, problem statement, research questions,
% contributions, scope and limitations, thesis roadmap. End with a per-chapter
% outline (the funnel: broad -> specific -> contributions -> roadmap).

% ---- Chapter 2: Literature Review / Related Work ----
\chapter{Literature Review}
\label{ch:related}
% TODO: a critical synthesis of prior work, organised by sub-problem.
% Identify the gap that this thesis addresses. Do NOT just list papers.

% ---- Chapter 3: Methods ----
\chapter{Methods}
\label{ch:methods}
% TODO: problem formulation, notation, proposed method, theoretical analysis,
% implementation details. State assumptions explicitly.

\section{Problem formulation}
% TODO.

\section{Proposed method}
% TODO.

\begin{algorithm}[t]
\caption{Outline of the proposed method}
\label{alg:method}
\begin{algorithmic}[1]
\State \textbf{Input:} dataset $\mathcal{D}$, hyper-parameters $\Theta$
\State initialise $\theta$
\While{not converged}
    \State $\theta \leftarrow \theta - \eta\,\nabla\mathcal{L}(\theta;\mathcal{D})$
\EndWhile
\State \textbf{return} $\theta$
\end{algorithmic}
\end{algorithm}

% ---- Chapter 4: Results / Experiments ----
\chapter{Results}
\label{ch:results}
% TODO: main results, ablation, qualitative examples, statistical significance,
% and a discussion of how the results answer each research question.

\begin{figure}[t]
\centering
% \includegraphics[width=0.8\textwidth]{figures/result.pdf}
\fbox{\rule{0pt}{6cm}\rule{0.7\textwidth}{0pt}}
\caption{Figure caption goes BELOW the figure. In a thesis, full-text-width
figures are common; use \texttt{\textbackslash textwidth}.}
\label{fig:result}
\end{figure}

\begin{table}[t]
\caption{Table caption goes ABOVE the table.}
\label{tab:results}
\centering
\begin{tabular}{lcc}
\toprule
Method   & Metric A       & Metric B       \\
\midrule
Baseline & 0.812          & 0.754          \\
Ours     & \textbf{0.891} & \textbf{0.832} \\
\bottomrule
\end{tabular}
\end{table}

% ---- Chapter 5: Conclusion ----
\chapter{Conclusion}
\label{ch:conclusion}
% TODO: recap the contributions and findings; state limitations; propose
% future work. A thesis conclusion is typically 3-6 pages.

% ==============================================================================
% Back matter
% ==============================================================================
\backmatter

% ---- Bibliography ----
\bibliography{references}
% Inline fallback (uncomment if not using BibTeX):
% \begin{thebibliography}{99}
% \bibitem{ref_example} F.~Author, \emph{Title of the paper}, Journal Name
%   \textbf{XX} (20YY) 1--10.
% \end{thebibliography}

% ---- Appendices ----
\appendix
\chapter{Supplementary Derivations}
\label{app:derivations}
% TODO: proofs, extended tables, dataset descriptions.

\chapter{Source Code Listings}
\label{app:code}
% TODO: key source-code listings; use the `listings` or `minted` package.

\end{document}
